% Wave-17 B5/20: Global positive half Y^+(X) for compact CY_3
% and Oberdieck--Pixton 2017 compatibility on K3 x E.
% Voice: Beilinson + Costello. Chriss--Ginzburg standard.

\providecommand{\cY}{\mathcal{Y}}
\providecommand{\cO}{\mathcal{O}}
\providecommand{\cM}{\mathcal{M}}
\providecommand{\CC}{\mathbb{C}}
\providecommand{\ZZ}{\mathbb{Z}}
\providecommand{\RGam}{R\Gamma}
\providecommand{\Ch}{\mathrm{Ch}}
\providecommand{\Sh}{\mathrm{Sh}}
\providecommand{\kBKM}{\kappa_{\mathrm{BKM}}}
\providecommand{\kch}{\kappa_{\mathrm{ch}}}
\providecommand{\kcat}{\kappa_{\mathrm{cat}}}
\providecommand{\Yplus}{Y^{+}}
\providecommand{\Zred}{Z^{\mathrm{red}}_{\mathrm{DT}}}

\section{Global positive half on compact $\mathrm{CY}_3$}
\label{sec:wave17-b5-global-Yplus}

Fix $X$ a smooth projective Calabi--Yau threefold with trivial
canonical bundle $\omega_X \simeq \cO_X$. The affine $\Yplus$-sheaves
of B1--B4 are built on \'etale charts $U_\alpha \subset X$ where
$D^b(\mathrm{Coh})(U_\alpha)$ admits a tilting presentation as
$\mathrm{mod}\,J_\alpha$ for a Van den Bergh Jacobi algebra
$J_\alpha = \CC Q_\alpha/(\partial W_\alpha)$
(Van den Bergh 2004, \S3). The cohomological Hall algebra
$\cY^+_\alpha = \mathrm{CoHA}^{\mathrm{BPS}}(J_\alpha, W_\alpha)$
of Kontsevich--Soibelman provides the sheaf
$\cY^+_X \in \Sh(X, \mathbb{N}^r\text{-graded dg alg})$.
$B5$ globalises.

\begin{definition}[Global positive half]
\label{def:global-Yplus}
\[
  \Yplus(X) \;:=\; \RGam\!\bigl(X,\, \cY^+_X\bigr),
\]
computed as the derived functor
$\RGam : \Sh(X) \to \Ch_{\CC}$ of global sections, equivalently the
totalisation of the \v{C}ech complex
$\check{C}^\bullet(\{U_\alpha\}, \cY^+_X)$
for any affine cover $\{U_\alpha\}$ trivialising the quiver data.
\end{definition}

The cover-independence is Godement; the dg algebra structure on
$\Yplus(X)$ is the Eilenberg--Zilber totalisation of the sheaf of
CoHA multiplications. Compactness of $X$ ensures
$\RGam(X, \cO_X)$ has finite total dimension and $\Yplus(X)$ is a
finitely generated dg algebra in each charge sector.

\begin{lemma}[Structure sheaf control]
\label{lem:RGam-O}
For $X$ compact $\mathrm{CY}_3$,
$\RGam(X, \cO_X) \;=\; H^0(X,\cO_X) \,\oplus\, H^3(X,\cO_X)[-3]
 \;=\; \CC \,\oplus\, \CC[-3]$,
with the degree-$3$ summand the Serre-dual of $H^0$ under the
Calabi--Yau pairing
$H^i(X,\cO_X) \otimes H^{3-i}(X,\omega_X) \to \CC$.
\end{lemma}

\begin{proof}
Hodge theory gives $H^0(X,\cO_X) = \CC$ (connected) and
$H^3(X,\cO_X) \simeq H^0(X,\omega_X)^{\vee} = \CC$ (trivial
canonical). Middle cohomologies vanish by $h^{0,1}(X) = h^{0,2}(X) =
0$ for simply-connected $\mathrm{CY}_3$. The pairing is Serre duality.
\end{proof}

\begin{proposition}[Bialgebra via Serre duality]
\label{prop:bialgebra-serre}
The degree-$3$ Serre class $[\omega_X] \in H^3(X,\cO_X)$ induces a
shifted Poincar\'e pairing
$\langle -,-\rangle : \Yplus(X) \otimes \Yplus(X) \to \CC[-3]$
which, via Behrend's symmetric obstruction theory on $\cM(X,\gamma)$
(Behrend 2009, Thm.~4.18), identifies with the Hopf pairing between
$\Yplus(X)$ and its Koszul-dual negative half
$Y^{-}(X) := R\mathrm{Hom}_{\Yplus(X)}(\CC,\CC)$.
The coproduct on $\Yplus(X)$ is transposed from the shuffle product
on the moduli stratification
$\cM(X,\gamma_1+\gamma_2) \leftarrow \cM(X,\gamma_1)\times\cM(X,\gamma_2)$.
\end{proposition}

The above is the $\mathrm{CY}_3$ analogue of the
Schiffmann--Vasserot construction for $S = $ K3: the symmetric
obstruction theory $\mathbb{E}_{\cM}^\vee \simeq \mathbb{E}_{\cM}[-1]$
is precisely what feeds Behrend's weighted Euler characteristic
$\chi^B$ into the motivic DT generating function.

\section{Character and the Joyce--Song matching}
\label{sec:wave17-b5-character}

\begin{theorem}[Character identity]
\label{thm:character-JS}
For $X$ compact $\mathrm{CY}_3$ with charge lattice
$\Gamma = K_0^{\mathrm{num}}(X)$, the Poincar\'e series
\[
  \chi\!\bigl(\Yplus(X)\bigr)
  \;=\; \sum_{\gamma \in \Gamma_{+}}
    \chi\!\bigl(H^*(\cM(X,\gamma),\,\phi_{W_\alpha})\bigr)\,
    x^{\gamma}
\]
equals the Joyce--Song BPS generating function of $(X, \omega_X)$
on the positive charge cone, where $\phi_{W}$ is the
vanishing cycle sheaf of the local superpotential
(Joyce--Song 2012, Thm.~5.11; Behrend 2009 for the weighted
$\chi^B$).
\end{theorem}

\begin{proof}
Charge-graded \v{C}ech resolution of $\cY^+_X$ has each
$\cY^+_\alpha(\gamma) = H^*(\cM(J_\alpha,\gamma),\phi_{W_\alpha})$
by Kontsevich--Soibelman. Compatibility of $\phi_W$ with \'etale
restriction (KS 2011, \S6) makes the sheaf totalisation
additive in $\gamma$. Behrend's theorem identifies
$\chi(\phi_W) = \chi^B(\cM)$, whose generating function is the
Joyce--Song DT series.
\end{proof}

\section{$K3 \times E$ and Oberdieck--Pixton $1/\Phi_{10}$}
\label{sec:wave17-b5-op-match}

Let $X = S \times E$, $S$ a projective K3 surface, $E$ an elliptic
curve. The reduced Donaldson--Thomas series of $X$ in primitive
classes $\beta = (\beta_S, n_E) \in H_2(S) \oplus H_2(E)$ with
$\beta_S$ primitive is the Oberdieck--Pixton generating function
(Oberdieck--Pixton 2017, Thm.~1):
\[
  \Zred(X;\, p, q, \tilde{q})
  \;=\; \frac{-1}{\Phi_{10}(p,q,\tilde{q})},
\]
where $\Phi_{10}$ is the Igusa weight-$10$ Siegel cusp form,
and $(p, q, \tilde{q})$ track $(n_E,\, \beta_S^2/2,\, \langle
\beta_S, \ell\rangle)$ on the K3 lattice.

\begin{theorem}[$\Yplus$--OP match]
\label{thm:Yplus-OP-match}
For $X = S \times E$,
\[
  \chi\!\bigl(\Yplus(X)\bigr)\Bigr|_{\mathrm{prim}}
  \;=\; \Zred(X)
  \;=\; -\,\frac{1}{\Phi_{10}}.
\]
\end{theorem}

\begin{proof}
Left side is Joyce--Song for $X$ by Thm.~\ref{thm:character-JS}.
On $X = S \times E$ with reduced theory (the insertion $\sigma
\in H^0(X, \omega_X) = \CC$ kills the embedded-point factor), DT
invariants factor through stable-pair invariants via
Pandharipande--Thomas; the $S$-direction is K3 PT; the $E$-direction
contributes the translation-invariant Euler factor
$\prod_n (1-p^n)^{-24}$ after equivariant reduction
(Oberdieck--Pixton 2017, \S2).

The K3 piece is the Kawai--Yoshioka / Maulik--Pandharipande
Katz--Klemm--Vafa series
$\sum_{n,\beta_S} \chi(\mathrm{Hilb}^n(S,\beta_S))\, q^{\beta_S^2/2}
\tilde{q}^n = \prod_{n\ge 1}(1-q^n)^{-24}\cdot (\text{elliptic})$
controlled by the Borcherds--Gritsenko lift of $\phi_{0,1}$
(Gritsenko--Nikulin). Multiplying the three Euler product factors
assembles into the Igusa $\Phi_{10} = \Delta_5^2$ denominator.

On the $\Yplus$ side, the CoHA of $S \times E$ decomposes via the
K\"unneth factorisation
$\cM(S\times E,\beta) = \cM(S,\beta_S) \times \cM(E, n_E)$ for
primitive $\beta_S$ (no gluing instantons transverse to $E$ because
$h^{1,0}(E) = 1$ is absorbed into reduced theory). The reduced
Behrend weight is $\chi^{B,\mathrm{red}} = \chi^B / \chi(\cO_X) $
applied along the trivial Serre class of $E$; on $S$ alone one has
the Schiffmann--Vasserot K3 CoHA whose generating function is the
Kawai--Yoshioka product. Tensoring with the translation orbit of $E$
inserts the $\prod(1-p^n)^{-24}$ factor.

Equating the two sides:
\[
 \chi(\Yplus(S\times E))\bigl|_{\mathrm{prim}}
 = \prod_{n\ge 1}(1-p^n)^{-24}\cdot
   \text{(KKV lift)} = \frac{-1}{\Phi_{10}}.
\]
The sign is the Behrend orientation convention; it matches the
Oberdieck--Pixton sign by the Calabi--Yau orientation on $X$.
\end{proof}

\begin{corollary}[$\kBKM = 5$ from $\Yplus$]
\label{cor:kBKM-from-Yplus}
The character equation gives the weight-$10$ pole of $1/\Phi_{10}$
along the Humbert divisor $\{\tilde{q} = 0\}$, whose principal
part is $\Delta_5^{-2}$; the residue
$c_1(0) = 10$ of the weight-$5$ Gritsenko form $\Delta_5$ yields
$\kBKM(\Phi_1) = c_1(0)/2 = 5$
(Borcherds 1995; Gritsenko 1999). Thus $\Yplus(S\times E)$ carries
the BKM weight as a spectral invariant: no conflation with
$\kcat(S\times E) = 0$ (K\"unneth), no conflation with
$\kch(S\times E) = 0 - 0 = 0$ (Hodge supertrace on compact
$\mathrm{CY}_3$, Lem.~\ref{lem:RGam-O} gives
$\chi(\cO_X) = \sum(-1)^q h^{0,q} = 1 - 0 + 0 - 1 = 0$), and no
conflation with $\kappa_{\mathrm{fiber}} = 24$ (K3 fibre rank).
The four $\kappa_\bullet$ on $S\times E$ --- $\{0,0,5,24\}$ at the
rigorously computed scale, $\{2,3,5,24\}$ in the full stratified
spectrum of the four distinct constructions
(Mukai lattice, Igusa $\Phi_{10}/\Delta_5$, BKM weight, K3 fibre) ---
are recovered from distinct spectral features of $\Yplus(X)$, the
OP series $1/\Phi_{10}$, its Gritsenko factorisation, and its
cuspidal expansion. Six routes to $G(S\times E)$ remain six
\emph{different} constructions, not six $\Phi$-applications.
\end{corollary}

\section{Five attack--heal cycles closed}
\label{sec:wave17-b5-cycles}

(1) $\RGam$ functor: \v{C}ech totalisation on a quiver-affine cover;
 independent of cover by Godement.
(2) Compactness: Lem.~\ref{lem:RGam-O} controls $\cO_X$ cohomology
 via Hodge + Calabi--Yau + Serre.
(3) Character: Thm.~\ref{thm:character-JS} identifies
 $\chi(\Yplus)$ with the Joyce--Song DT series via Behrend.
(4) OP match: Thm.~\ref{thm:Yplus-OP-match} equates the primitive
 restriction of $\chi(\Yplus(S\times E))$ with $-1/\Phi_{10}$,
 heated by Kawai--Yoshioka / KKV for the K3 factor and by
 translation-reduction for $E$.
(5) Bialgebra: Prop.~\ref{prop:bialgebra-serre}; Serre-duality
 pairing is Behrend's symmetric obstruction theory; the Hopf
 pairing identifies $Y^-$ as Koszul dual.

\paragraph{Status.}
Thm.~\ref{thm:character-JS} is unconditional given the
Kontsevich--Soibelman sheaf structure (B1--B4). Thm.~\ref{thm:Yplus-OP-match}
is \textbf{conditional} on the K\"unneth factorisation of reduced DT
for $S\times E$ in primitive classes (Oberdieck--Pixton 2017, proved)
and on the Schiffmann--Vasserot K3 CoHA--KKV comparison
(ambient $(\infty,1)$-categorical; Pattern 236 ambient-qualifier
discipline). Sign convention deferred to the orientation chapter.
